\documentclass[svenska]{style/LTHthesis}

% Preamble
% from template
\usepackage[T1]{fontenc}%
\usepackage{mathptmx, helvet}%
\usepackage{graphicx}%

% To include all authors in \fullcite command
\preto\fullcite{\AtNextCite{\defcounter{maxnames}{99}}}%

%%%%%%%%%%%%%%%%%%%%%%%%%
%%% packages from:
%%%    kappa/
%%% as you add the papers you'll have to adjust them
%%%%%%%%%%%%%%%%%%%%%%%%%

\usepackage{fontawesome5}%
\usepackage{verbatimbox}%

%%%%%%%%%%%%%%%%%%%%%%%%%
%%% packages copy pasted from:
%%%    ecrts21/og/pkg.tex
%%% as you add the papers you'll have to adjust them
%%%%%%%%%%%%%%%%%%%%%%%%%

\usepackage{amssymb, mathtools}%
\usepackage[inline]{enumitem}%      %inline enumeration
\usepackage[table]{xcolor}%         %The 'table' option for 'xcolor' loads the 'colortbl' package, so there's no need to load it separately.
\usepackage{xspace}%
\usepackage{afterpage}%

%%%%%%%%%%%%%%%%%%%%%%%%%
%%% packages copy pasted from:
%%%    rtas22a/og/pkg.tex
%%% as you add the papers you'll have to adjust them
%%%%%%%%%%%%%%%%%%%%%%%%%

\usepackage{xfrac}%                 % For nice fractions
%\usepackage{colortbl} % For colourful tables % NOTE: Not needed after the table option was introduced for xcolor

%%%%%%%%%%%%%%%%%%%%%%%%%
%%% packages copy pasted from:
%%%    rtas22b/og/pkg.tex
%%% as you add the papers you'll have to adjust them
%%%%%%%%%%%%%%%%%%%%%%%%%

\usepackage{algorithm}%             % Algorithm related
\usepackage[noend]{algpseudocode}%  % Algorithm related

%%%%%%%%%%%%%%%%%%%%%%%%%
%%% packages copy pasted from:
%%%    lcss22/og/pkg.tex
%%% as you add the papers you'll have to adjust them
%%%%%%%%%%%%%%%%%%%%%%%%%

\usepackage{lscape}%

%%%%%%%%%%%%%%%%%%%%%%%%%
%%% Tikz definitions and libraries
%%%%%%%%%%%%%%%%%%%%%%%%%

% TIKZ settings
\usepackage{tikz, pgfplots, pgfplotstable}
\usepackage{shellesc} %As suggested in the wiki

% externalization
\usetikzlibrary{external}
\tikzexternalize[prefix=externalised/]

%%%%%%%%%%%%%%%%%%%%%%%%%
%%% packages copy pasted from:
%%%    ecrts21/og/pkg.tex
%%% as you add the papers you'll have to adjust them
%%%%%%%%%%%%%%%%%%%%%%%%%

%%%%%%%%%%%%%%%%%%%%%%%%%%
% libraries
\pgfplotsset{compat=1.16}
\usetikzlibrary{shapes, fit}
\usetikzlibrary{arrows, arrows.meta}
\usetikzlibrary{decorations.markings}
\usepgfplotslibrary{groupplots}

% Extracts the max value from the input and marks it in the plot
\pgfplotsset{
  mark max/.style={
    scatter/@pre marker code/.code={%
    \ifx\pgfplotspointmeta\pgfplots@metamax
        \def\markopts{}%
        \node [anchor=south, xshift=1mm] { \pgfmathprintnumber[precision=1, fixed zerofill, ultra thick]{\pgfplotspointmeta} };
    \else
        \def\markopts{mark=none}
    \fi
        \expandafter\scope\expandafter[\markopts]
    },%
    scatter/@post marker code/.code={ \endscope },
    scatter
  }
}

% Read element from specific place in data table
\newcommand*{\ReadOutElement}[4]{%
    \pgfplotstablegetelem{#2}{[index]#3}\of{#1}%
    \let#4\pgfplotsretval
}

%%%%%%%%%%%%%%%%%%%%%%%%%
%%% packages copy pasted from:
%%%    rtas22a/og/pkg.tex
%%% as you add the papers you'll have to adjust them
%%%%%%%%%%%%%%%%%%%%%%%%%

\usetikzlibrary{shapes.misc,backgrounds}
%

%%%%%%%%%%%%%%%%%%%%%%%%%
%%% Theorem formats
%%%%%%%%%%%%%%%%%%%%%%%%%

\newtheorem{theorem}{Theorem}
\newtheorem{corollary}{Corollary}
\newtheorem{problem}{Problem}
\newtheorem{lemma}{Lemma}
\newtheorem{remark}{Remark}
\newtheorem{definition}{Definition}
%
%
%%%%%%%%%%%%%%%%%%%%%%%%%
%%% Unified notation
%%%%%%%%%%%%%%%%%%%%%%%%%

%% Mathematics
% Automaton related
\newcommand{\GG}[1]{\mathcal{G}_{#1}} % Automaton
\newcommand{\VV}[1]{V_{#1}} % Nodes
\newcommand{\EE}[1]{E_{#1}} % Edges

% Language related
\newcommand{\cM}{\texttt{M}}
\newcommand{\cH}{\texttt{H}}
\newcommand{\cR}{\texttt{R}}
\newcommand{\cT}{\texttt{T}}
\newcommand{\cX}{\texttt{X}}


% Shorthand
\newcommand{\R}{\mathbb{R}}
\newcommand{\N}{\mathbb{N}}
\newcommand{\T}{\mathrm{\small{T}}}
\newcommand{\nequiv}{\not\equiv}  

% Functions
\newcommand{\LL}{\mathcal{L}_2}
\newcommand{\funof}[1]{\left( #1 \right)}
\newcommand{\abs}[1]{\left|#1\right|}
\newcommand{\norm}[1]{\lVert#1\rVert}
\newcommand{\E}[1]{\mathbb{E}\left[ #1 \right]}
\newcommand{\trace}[1]{\mathrm{tr}\left(#1\right)}
\newcommand{\diag}[1]{\mathrm{diag}\left\{#1\right\}}
\DeclareMathOperator*{\argmax}{argmax}
\DeclareMathOperator*{\argmin}{argmin}
\newcommand{\ourmod}[1]{\ \mathrm{mod}\ #1}
\DeclarePairedDelimiter\ceil{\lceil}{\rceil}
\DeclarePairedDelimiter\floor{\lfloor}{\rfloor}
\newcommand{\overbar}[1]{\mkern 1.5mu\overline{\mkern-1.5mu#1\mkern-1.5mu}\mkern 1.5mu}
\newcommand{\fact}[1]{#1!}

%% Control
\newcommand{\plant}{\mathcal{P}}
\newcommand{\ctrler}{\mathcal{C}}

%% Code
\newcommand{\code}[1]{\texttt{\textbf{#1}}} % code environments
\newcommand{\tool}{\code{WeaklyHard.jl}}

%% Real-Time
% Strategies
\newcommand{\strat}{\mathcal{H}} 

% Constraint names
\newcommand{\tAH}{\code{AnyHit}}
\newcommand{\tAM}{\code{AnyMiss}}
\newcommand{\tRH}{\code{RowHit}}
\newcommand{\tRM}{\code{RowMiss}}

% Constraints related 
\newcommand{\anyhit}[1]{\binom{x_{#1}}{k_{#1}}} 
\newcommand{\anymiss}[1]{\overbar{\binom{x_{#1}}{k_{#1}}}}
\newcommand{\rowhit}[1]{\genfrac{<}{>}{0pt}{}{x_{#1}}{k_{#1}}}
\newcommand{\rowmiss}[1]{\overbar{\genfrac{<}{>}{0pt}{}{x_{#1}}{k_{#1}}}}
\newcommand{\eanyhit}[2]{\binom{x_{#1}}{k_{#1}}^{#2}}
\newcommand{\eanymiss}[2]{\overbar{\binom{x_{#1}}{k_{#1}}}^{#2}}
\newcommand{\erowhit}[2]{\genfrac{<}{>}{0pt}{}{x_{#1}}{k_{#1}}^{#2}}
\newcommand{\erowmiss}[2]{\overbar{\genfrac{<}{>}{0pt}{}{x_{#1}}{k_{#1}}}^{#2}}

\newcommand{\lweak}{\underline{\lambda}} 
\newcommand{\lhard}{\overbar{\lambda}}

\newcommand{\sset}[2]{\mathcal{S}_{#1}\funof{#2}}

%% Blame commands
\newcommand{\fix}[1]{\textcolor{red}{\textbf{FIX! #1}}} 
\newcommand{\nv}[1]{\textcolor{green!60!blue}{\textbf{NV:} #1}}
%

% Document
\begin{document}

\noindent
\begin{minipage}[]{0.25\textwidth}
    \noindent\includegraphics[width=20mm]{style/LUCsv}
    \par\vspace*{3mm} 
\end{minipage}\hfill
\begin{minipage}[]{0.74\textwidth}
    \vspace*{-3mm}
    \sffamily
    \begin{flushleft}
        \huge Vad är det värsta som kan hända när datorn krånglar? \\[2mm]
        \Large Nils Vreman\\[1mm]
        \normalsize Institutionen för Reglerteknik
    \end{flushleft}
\end{minipage}

\par\vspace{1mm}  

\begin{sffamily}\noindent
    Populärvetenskaplig sammanfattning av doktorsavhandlingen \emph{Analysis of Embedded Controllers Subject to Computational Overruns}, juni 2023.
    Avhandlingen kan laddas ner från: \url{http://www.control.lth.se/publications}
\end{sffamily}

\vspace{2mm}\hrule\vspace{3mm}

\noindent
Datorer i olika storlekar och med olika styrka finns runt omkring oss hela tiden.
Allt från mobilen i din ficka till bilen du åker med till jobbet innehåller många små datorer, också kallade \emph{inbyggda system}.
Varje litet inbyggt system har uppgifter som måste uträttas, t.ex., beräkna hastigheten.
I de flesta fall finns det krav och specifikationer på att beräkningarna måste uträttas inom en viss tidsram.
Till exempel, om bromsen trycks in förväntas bilen sakta ner direkt.
I verkligheten tar det lite tid för det inbyggda systemet att beräkna nedbromsningshastigheten innan kommandot kan skickas till bromssystemet.
Systemet är dock alltid designat med hänsyn till den extra tiden det tar att uträtta hela uppgiften.

\begin{wrapfigure}{r}{0.45\textwidth}
    \vspace*{-2mm}
    \begin{centering}
        \includegraphics[width=0.45\textwidth]{figs/microcontroller.jpg}
    \end{centering}\small\par\medskip
        \noindent\emph{Exempel på ett inbyggt system. På bilden ses en Raspberry Pi Pico där det svarta chippet i mitten (storlek $7\times 7$~mm) är en mikrokontroller, dvs., ett inbyggt system.}\footnotemark
    \vspace*{-2mm}
\end{wrapfigure}
\footnotetext{Laserlicht, CC BY-SA 4.0, \url{https://creativecommons.org/licenses/by-sa/4.0/}, via Wikimedia Commons}

Tyvärr kan inbyggda system krångla, precis som en persondator.
Bland annat kan bakgrundsprogram, virus, och föråldrad hårdvara/programvara störa beräkningen.
Det kan försena beräkningen, och i värsta fall medföra att uppgiften inte blir klar alls.
Vad händer då om det inbyggda systemet krånglar på ett sådant sätt att tidsramen för beräkningen inte respekteras?
Just den frågan besvaras i den här avhandlingen.

Många inbyggda system sitter i apparater där det inte gör så mycket om något krånglar.
Till exempel har mycket hemelektronik (så som moderna stereosystem, TV-apparater, och kylskåp) inbyggda system för att styra funktionaliteten och kommunicera med andra smarta system.
Det är irriterande om smart-TVns bild blir oskarp, men det är inget som påverkar oss märkbart.
Värre konsekvenser uppstår om det inbyggda systemet krånglar i bilens bromssystem, en pacemaker, eller i ett flygplans styrsystem.
Då kan det uppstå verklig fara för människor i närheten och värstascenariot inkluderar även dödsfall.

Den här avhandlingen utvecklar matematiska verktyg och metoder för att förbättra analysen av dessa säkerhetskritiska inbyggda system.
Det är livsviktigt att ingen person kommer till skada och därför fokuserar analysen på att se till att hela systemet är säkert, även under väldigt dåliga förhållanden.
Dåliga förhållanden innebär här, bland annat, att det inbyggda systemet inte lyckas beräkna något under en väldigt lång tid på grund av, t.ex., en hacker-attack.
Med förhandsinformation om hur länge systemet är säkert när något går fel är det möjligt att utveckla metoder för att säkert stänga av eller starta om systemet innan någon kommer till skada.

Ett säkerhetskritiskt system blir inte nödvändigtvis osäkert bara för att beräkningen tar lite extra tid då och då, men det kan göra att upplevelsen blir väldigt obehaglig för användaren.
Till exempel kan bilens nedbromsning bli ryckig och oregelbunden.
På lång sikt är det då troligt att systemet bryts ner och går sönder på grund av utförandets oregelbundenhet, och som konsument överger du antagligen tillverkaren.
Detta kan kosta tillverkaren väldigt mycket pengar och försämrar generellt sett företagets rykte.
Därför räcker det inte att enbart se till att produkten är helt säker, den måste också nå upp till en viss kvalitet för att det ska bli lönsamt att producera den.

I avhandlingen utformas både algoritmer för att förbättra produktkvaliteten och metoder för att kvalitetsgranska slutprodukten.
Genom att förändra beräkningarna som utförs av det inbyggda systemet i realtid visas att det är möjligt att kompensera för den försämrade prestandan.
Nya metoder för att undersöka prestandan utvecklas och används för att garantera att algoritmerna fungerar enligt förväntan när det inbyggda systemet krånglar.
För att bekräfta att avhandlingens bidrag går att använda i verkligheten testas både algoritmerna, säkerhetsanalysen och prestandaanalysen på flertalet verkliga system.
Experimenten visar att metoderna både är enkla att använda och ger en bra bild av hur det verkliga systemet påverkas av problem med de inbyggda systemen.

\end{document}
